% !TEX root = main.tex
\section{提案：最適化フレームワークへのフィードフォワード型構成の導入}
\label{sec:proposal}

\subsection{既存の最適化フレームワーク}

本研究室では，PLLのジッタと消費電力の間のトレードオフを定量的に比較するための最適化フレームワークが提案されている\cite{zhu_moea_2025}．このフレームワークは，多目的進化アルゴリズム（Multi-Objective Evolutionary Algorithm, MOEA，具体的にはNSGA-II）を用いて，PLLの各ブロックの消費電力やループ帯域周波数などを決定変数として探索し，積分ジッタとトータル消費電力の2つの目的関数について，一方を改善すればもう一方が悪化するという意味で最適な解の集合（パレートフロント）を求めるものである．これにより，従来は設計者が数式変形と勘に頼って行っていたジッタ・電力トレードオフの評価を，体系的かつ網羅的に行うことができる．

このフレームワークは，従来型の整数分周比PLL，DTCベースPLL，およびHMベースのDual-Feedback PLLを対象に，ジッタ・電力のパレートフロントを比較しており，各アーキテクチャがどのような条件で優位に立つかを議論している．しかし，第\ref{sec:methods}節で述べたフィードフォワード型構成は，このフレームワークにはまだ組み込まれていない．

\subsection{本研究で行うこと}

そこで本研究では，この最適化フレームワークに，フィードフォワード型構成のノイズ伝達関数を新たに組み込む．具体的には，補助PLLの雑音源から出力までの伝達関数に，フィードフォワードによるキャンセル項（ゲイン整合条件を含む）を追加したモデルを構築し，これを目的関数（積分ジッタ・消費電力）の計算に反映させる．

これにより，(1)DTCベースPLL，(2)フィードフォワードなしのHMベースDual-Feedback PLL，(3)フィードフォワード型構成を備えたDual-Feedback PLL，の3アーキテクチャについて，同一の枠組みの上でジッタ・電力パレートフロントを比較することが可能になる．この比較を通じて，フィードフォワード型構成がどのジッタ・電力の領域で他の構成に対して優位に立つのかを明らかにすることが，本研究の中心的な取り組みである．
